%%%%%%%%%%%%%%%%%%%%%%%%%%%%%%%%%%%%%%%%%%%%%%%%%%%%%%%%%%%%%%%%%%%%%%%%%%%%%%%%
%% Capitulo 5: Documento de Design Detalhado (SDD)
%% Sistema: PostoSmart - Automação e Gestão de Postos de Combustíveis
%% Autor: Diogo Santos Pires Jandiroba
%%%%%%%%%%%%%%%%%%%%%%%%%%%%%%%%%%%%%%%%%%%%%%%%%%%%%%%%%%%%%%%%%%%%%%%%%%%%%%%%

\chapter{Documento de Design Detalhado (SDD)}
\label{cap:design_tecnico}

\section{Diagrama de Classes em Nível de Projeto}

A \autoref{fig:classes_projeto_posto} apresenta a estrutura de classes de implementação do módulo de automação de pista e faturamento, detalhando visibilidades, tipos e assinaturas de métodos.

% PLACEHOLDER STARUML v7.0: Diagrama de Classes em Nível de Projeto
% Ferramenta: StarUML v7.0 -> Model -> Add Diagram -> Class Diagram
% Arquivo: Imagens/cls_projeto_postosmart.png (Exportar em PNG 300 DPI ou PDF vetorial)
% Descomente a linha abaixo apos salvar o arquivo no diretorio Imagens/:
% \incluirdiagrama{Imagens/cls_projeto_postosmart.png}{Diagrama de Classes em Nível de Projeto}{fig:classes_projeto_posto}
\begin{figure}[htbp]
    \centering
    \fbox{\parbox{0.9\textwidth}{\centering\vspace{1cm}%
        \textbf{[Diagrama de Classes em Nível de Projeto -- StarUML v7.0]}\\%
        \textit{Modelar no StarUML v7.0 e exportar para: \texttt{Imagens/cls_projeto_postosmart.png}}\\%
        \small Menu: \texttt{File -> Export Diagram as -> PNG...} (Fundo branco, alta resolução 300 DPI)\\%
        \vspace{0.3cm}%
        \footnotesize Para ativar a imagem real, descomente no arquivo \texttt{.tex}:\\
        \texttt{\textbackslash incluirdiagrama\{Imagens/cls_projeto_postosmart.png\}\{Diagrama de Classes em Nível de Projeto\}\{fig:classes_projeto_posto\}}%
    \vspace{1cm}}}
    \caption{Diagrama de Classes em Nível de Projeto (Placeholder StarUML v7.0)}
    \label{fig:classes_projeto_posto}
\end{figure}

\section{Diagrama de Sequência: Autorização e Captura de Abastecimento}

A \autoref{fig:sequencia_abastecimento_posto} demonstra a dinâmica temporal de troca de mensagens entre o hardware de pista, o concentrador e o PDV.

% PLACEHOLDER STARUML v7.0: Diagrama de Sequência da Autorização e Captura
% Ferramenta: StarUML v7.0 -> Model -> Add Diagram -> Sequence Diagram
% Arquivo: Imagens/seq_abastecimento.png (Exportar em PNG 300 DPI ou PDF vetorial)
% Descomente a linha abaixo apos salvar o arquivo no diretorio Imagens/:
% \incluirdiagrama{Imagens/seq_abastecimento.png}{Diagrama de Sequência da Autorização e Captura de Abastecimento}{fig:sequencia_abastecimento_posto}
\begin{figure}[htbp]
    \centering
    \fbox{\parbox{0.9\textwidth}{\centering\vspace{1cm}%
        \textbf{[Diagrama de Sequência: Autorização e Captura -- StarUML v7.0]}\\%
        \textit{Modelar no StarUML v7.0 e exportar para: \texttt{Imagens/seq_abastecimento.png}}\\%
        \small Menu: \texttt{File -> Export Diagram as -> PNG...} (Fundo branco, alta resolução 300 DPI)\\%
        \vspace{0.3cm}%
        \footnotesize Para ativar a imagem real, descomente no arquivo \texttt{.tex}:\\
        \texttt{\textbackslash incluirdiagrama\{Imagens/seq_abastecimento.png\}\{Diagrama de Sequência da Autorização e Captura de Abastecimento\}\{fig:sequencia_abastecimento_posto\}}%
    \vspace{1cm}}}
    \caption{Diagrama de Sequência da Autorização e Captura de Abastecimento (Placeholder StarUML v7.0)}
    \label{fig:sequencia_abastecimento_posto}
\end{figure}

\section{Diagrama de Máquina de Estados: Ciclo de Vida do Bico}

A \autoref{fig:maquina_estados_bico} especifica os estados operacionais finitos de um bico de abastecimento.

% PLACEHOLDER STARUML v7.0: Diagrama de Máquina de Estados do Bico
% Ferramenta: StarUML v7.0 -> Model -> Add Diagram -> Statechart Diagram
% Arquivo: Imagens/dsm_bico_combustivel.png (Exportar em PNG 300 DPI ou PDF vetorial)
% Descomente a linha abaixo apos salvar o arquivo no diretorio Imagens/:
% \incluirdiagrama{Imagens/dsm_bico_combustivel.png}{Diagrama de Máquina de Estados do Bico de Combustível}{fig:maquina_estados_bico}
\begin{figure}[htbp]
    \centering
    \fbox{\parbox{0.9\textwidth}{\centering\vspace{1cm}%
        \textbf{[Diagrama de Máquina de Estados: Bico de Combustível -- StarUML v7.0]}\\%
        \textit{Modelar no StarUML v7.0 e exportar para: \texttt{Imagens/dsm_bico_combustivel.png}}\\%
        \small Menu: \texttt{File -> Export Diagram as -> PNG...} (Fundo branco, alta resolução 300 DPI)\\%
        \vspace{0.3cm}%
        \footnotesize Para ativar a imagem real, descomente no arquivo \texttt{.tex}:\\
        \texttt{\textbackslash incluirdiagrama\{Imagens/dsm_bico_combustivel.png\}\{Diagrama de Máquina de Estados do Bico de Combustível\}\{fig:maquina_estados_bico\}}%
    \vspace{1cm}}}
    \caption{Diagrama de Máquina de Estados do Bico de Combustível (Placeholder StarUML v7.0)}
    \label{fig:maquina_estados_bico}
\end{figure}

\section{Diagrama de Atividades: Pagamento e Emissão Fiscal}

A \autoref{fig:atividades_fechamento_posto} modela o fluxo de fechamento de venda no caixa com tratamento de contingência fiscal.

% PLACEHOLDER STARUML v7.0: Diagrama de Atividades de Fechamento de Venda
% Ferramenta: StarUML v7.0 -> Model -> Add Diagram -> Activity Diagram
% Arquivo: Imagens/act_fechamento_venda.png (Exportar em PNG 300 DPI ou PDF vetorial)
% Descomente a linha abaixo apos salvar o arquivo no diretorio Imagens/:
% \incluirdiagrama{Imagens/act_fechamento_venda.png}{Diagrama de Atividades do Fechamento de Venda e Emissão Fiscal}{fig:atividades_fechamento_posto}
\begin{figure}[htbp]
    \centering
    \fbox{\parbox{0.9\textwidth}{\centering\vspace{1cm}%
        \textbf{[Diagrama de Atividades: Fechamento de Venda e Fiscal -- StarUML v7.0]}\\%
        \textit{Modelar no StarUML v7.0 e exportar para: \texttt{Imagens/act_fechamento_venda.png}}\\%
        \small Menu: \texttt{File -> Export Diagram as -> PNG...} (Fundo branco, alta resolução 300 DPI)\\%
        \vspace{0.3cm}%
        \footnotesize Para ativar a imagem real, descomente no arquivo \texttt{.tex}:\\
        \texttt{\textbackslash incluirdiagrama\{Imagens/act_fechamento_venda.png\}\{Diagrama de Atividades do Fechamento de Venda e Emissão Fiscal\}\{fig:atividades_fechamento_posto\}}%
    \vspace{1cm}}}
    \caption{Diagrama de Atividades do Fechamento de Venda e Emissão Fiscal (Placeholder StarUML v7.0)}
    \label{fig:atividades_fechamento_posto}
\end{figure}

\section{Esquema Relacional e Mapeamento ORM}

A \autoref{fig:er_banco_posto} exibe o modelo relacional de tabelas físicas com restrições de integridade referencial.

% PLACEHOLDER STARUML v7.0: Diagrama Entidade-Relacionamento / DDL Relacional
% Ferramenta: StarUML v7.0 -> Model -> Add Diagram -> Class Diagram (ou ERD)
% Arquivo: Imagens/cls_er_relacional.png (Exportar em PNG 300 DPI ou PDF vetorial)
% Descomente a linha abaixo apos salvar o arquivo no diretorio Imagens/:
% \incluirdiagrama{Imagens/cls_er_relacional.png}{Diagrama Entidade-Relacionamento do Banco de Dados Relacional}{fig:er_banco_posto}
\begin{figure}[htbp]
    \centering
    \fbox{\parbox{0.9\textwidth}{\centering\vspace{1cm}%
        \textbf{[Diagrama ER / Modelo Físico Relacional -- StarUML v7.0]}\\%
        \textit{Modelar no StarUML v7.0 e exportar para: \texttt{Imagens/cls_er_relacional.png}}\\%
        \small Menu: \texttt{File -> Export Diagram as -> PNG...} (Fundo branco, alta resolução 300 DPI)\\%
        \vspace{0.3cm}%
        \footnotesize Para ativar a imagem real, descomente no arquivo \texttt{.tex}:\\
        \texttt{\textbackslash incluirdiagrama\{Imagens/cls_er_relacional.png\}\{Diagrama Entidade-Relacionamento do Banco de Dados Relacional\}\{fig:er_banco_posto\}}%
    \vspace{1cm}}}
    \caption{Diagrama Entidade-Relacionamento do Banco de Dados Relacional (Placeholder StarUML v7.0)}
    \label{fig:er_banco_posto}
\end{figure}

\paragraph{Script DDL SQL de Criação de Tabelas (PostgreSQL / SQLite):}
\begin{lstlisting}[language=SQL]
CREATE TABLE tb_bico (
    id_bico INT PRIMARY KEY,
    id_tanque UUID NOT NULL REFERENCES tb_tanque(id_tanque),
    numero_bico INT NOT NULL UNIQUE,
    preco_litro DECIMAL(10,2) NOT NULL,
    status VARCHAR(20) NOT NULL DEFAULT 'LIVRE'
);

CREATE TABLE tb_abastecimento (
    id_abastecimento UUID PRIMARY KEY,
    id_bico INT NOT NULL REFERENCES tb_bico(id_bico),
    volume_litros DECIMAL(10,3) NOT NULL,
    valor_total DECIMAL(10,2) NOT NULL,
    data_hora TIMESTAMP NOT NULL DEFAULT CURRENT_TIMESTAMP,
    status VARCHAR(20) NOT NULL DEFAULT 'PENDENTE'
);
\end{lstlisting}
